\chapter{Complexity — Circuit Eval Spec}
%%% generated by the blueprint pipeline from TCSlib.Complexity.CircuitComplexity.CircuitEvalSpec
\section{Overview}

Declarations of the Lean module \texttt{TCSlib.Complexity.CircuitComplexity.CircuitEvalSpec}.

\section{Declarations}

\begin{definition}[Evaluation phases]
\lean{BoolCircuit.CircuitEval.Ph}
\label{BoolCircuit.CircuitEval.Ph}
\leanok
\uses{BoolCircuit.GateKind}
Define the finite set of evaluation phases consisting of $n0$, $sep$, $skipN$, $g$,
$k1$, $k2(b)$ for a Boolean $b$, $a(k,c,acc)$ and $u(k,c,acc)$ for a gate kind $k$, a
three-valued index $c$, and a Boolean accumulator $acc$, together with $o$ and $fin(v)$
for a Boolean $v$. The gate kind ranges over AND, OR, and NOT.
\end{definition}

\begin{definition}[Machine evaluation states]
\lean{BoolCircuit.CircuitEval.St}
\label{BoolCircuit.CircuitEval.St}
\leanok
\uses{BoolCircuit.CircuitEval.Ph}
Define the finite set of machine states consisting of $rd1(\phi)$, $rd2(\phi,b)$,
$rewM$, $copy$, $rewV(d)$, $rewI0$, $rewI$, and $app(v)$, where $\phi$ is an evaluation
phase, $b$ and $v$ are Booleans, and $d$ is either absent or an evaluation phase.
\end{definition}

\begin{definition}[Gate kind from two bits]
\lean{BoolCircuit.CircuitEval.kindOf}
\label{BoolCircuit.CircuitEval.kindOf}
\leanok
\uses{BoolCircuit.GateKind}
For Boolean inputs $x$ and $y$, define the gate kind to be AND when
$(x,y)=(	ext{false},	ext{false})$, OR when $(x,y)=(	ext{false},	ext{true})$, and NOT
whenever $x=	ext{true}$.
\end{definition}

\begin{definition}[Gate accumulator initialization]
\lean{BoolCircuit.CircuitEval.accInit}
\label{BoolCircuit.CircuitEval.accInit}
\leanok
\uses{BoolCircuit.GateKind}
Initialize the accumulator for an OR gate to false, and for an AND or NOT gate to true.
\end{definition}

\begin{definition}[Gate accumulator update]
\lean{BoolCircuit.CircuitEval.accComb}
\label{BoolCircuit.CircuitEval.accComb}
\leanok
\uses{BoolCircuit.GateKind}
Given an accumulator value $a$ and an input bit $v$, update an OR gate’s accumulator to
$a \lor v$; for an AND or NOT gate, update it to $a \land v$.
\end{definition}

\begin{definition}[Gate accumulator finalization]
\lean{BoolCircuit.CircuitEval.accFin}
\label{BoolCircuit.CircuitEval.accFin}
\leanok
\uses{BoolCircuit.GateKind}
To finalize an accumulator value $a$, return its Boolean negation for a NOT gate and
return $a$ unchanged for an AND or OR gate.
\end{definition}

\begin{definition}[Machine action record]
\lean{BoolCircuit.CircuitEval.LA}
\label{BoolCircuit.CircuitEval.LA}
\leanok
\uses{BoolCircuit.CircuitEval.St}
An action record consists of an optional value-tape write, an optional mark-tape write,
a sign specifying the movement of both work heads, an optional emitted Boolean symbol,
and an optional next machine state.
\end{definition}

\begin{definition}[Rejecting machine action]
\lean{BoolCircuit.CircuitEval.LA.rej}
\label{BoolCircuit.CircuitEval.LA.rej}
\leanok
\uses{BoolCircuit.CircuitEval.LA}
The rejecting action writes nothing on either work tape, leaves both work heads
stationary, emits the Boolean symbol false, and specifies no next state.
\end{definition}

\begin{definition}[Continue with a machine state]
\lean{BoolCircuit.CircuitEval.LA.go}
\label{BoolCircuit.CircuitEval.LA.go}
\leanok
\uses{BoolCircuit.CircuitEval.LA, BoolCircuit.CircuitEval.St}
For any machine state $q$, the action that continues with state $q$ writes nothing on
either work tape, leaves both work heads stationary, emits no symbol, and specifies $q$
as the next state.
\end{definition}

\begin{definition}[Circuit-evaluation machine transition]
\lean{BoolCircuit.CircuitEval.lact}
\label{BoolCircuit.CircuitEval.lact}
\leanok
\uses{BoolCircuit.CircuitEval.LA, BoolCircuit.CircuitEval.LA.go, BoolCircuit.CircuitEval.LA.rej, BoolCircuit.CircuitEval.Ph, BoolCircuit.CircuitEval.accComb, BoolCircuit.CircuitEval.accFin, BoolCircuit.CircuitEval.accInit, BoolCircuit.CircuitEval.kindOf}
The transition is a function of a phase and three optional Boolean tape readings. In
phase $n0$, a true value writes true on the mark tape, moves both heads positively, and
enters $rd1(n0)$; a false value enters $rd1(sep)$. In phase $sep$, a present value
enters $rd1(sep)$, while an absent value moves both heads negatively and enters $rewM$.
In phase $skipN$, true continues in that phase and false enters $rd1(g)$; in phase $g$,
true enters $rd1(k1)$ and false enters $rd1(o)$. In phase $k1$, a present bit $b$ enters
$rd1(k2(b))$. In phase $k2$, with stored bit $b_1$ and current bit $b_2$, the transition
rejects if both are true; otherwise it enters $rd1(a(k,0,accInit(k)))$, where $k$ is the
gate kind determined by $(b_1,b_2)$. In phase $a(k,c,acc)$, a true value rejects if
$c=2$ and otherwise enters $rd1(u(k,c,acc))$; a false value rejects when $k$ is NOT and
$c\ne1$, and otherwise enters $app(accFin(k,acc))$. In phase $u(k,c,acc)$, a true value
together with a present mark reading moves both heads positively and continues in
$rd1(u(k,c,acc))$; a false value together with a present value reading $v$ and absent
third reading writes true on the mark tape, moves both heads negatively, and enters
$rewV(a(k,c+1,accComb(k,acc,v)))$. In phase $o$, a true value together with a present
mark reading moves both heads positively and continues in $rd1(o)$; a false value
together with a present value reading $v$ enters $rd1(fin(v))$. In phase $fin(v)$ with
no value reading, the machine emits $v$ and halts. Every input pattern not specified
above produces the rejecting action.
\end{definition}

\begin{definition}[Evaluator configuration]
\lean{BoolCircuit.CircuitEval.AS}
\label{BoolCircuit.CircuitEval.AS}
\leanok
\uses{BoolCircuit.CircuitEval.Ph}
An evaluator configuration consists of a phase, a finite list of Boolean values already
computed, a finite list of marked natural-numbered tape cells, and a natural-numbered
work-head position.
\end{definition}

\begin{definition}[Mark-tape lookup]
\lean{BoolCircuit.CircuitEval.markCell}
\label{BoolCircuit.CircuitEval.markCell}
\leanok
Given a list $M$ of marked natural-numbered cells and a position $j$, the lookup returns
true if $j$ belongs to $M$, and is undefined otherwise.
\end{definition}

\begin{definition}[Abstract evaluator update]
\lean{BoolCircuit.CircuitEval.absOf}
\label{BoolCircuit.CircuitEval.absOf}
\leanok
\uses{BoolCircuit.CircuitEval.AS, BoolCircuit.CircuitEval.LA}
Given an evaluator configuration and an action, the update either returns an output
Boolean or a new configuration. If the action has no next state, it outputs the action's
symbol, using false when that symbol is absent. A read action sets the phase to its
specified phase, optionally marks the current head cell, and advances the head by one
exactly when the head move is positive. A value-rewind action with a specified phase
performs the same optional marking and resets the head to zero. An append action appends
its Boolean value to the computed values, clears the marked cells, resets the head to
zero, and sets the phase to the gate-evaluation phase; every other action returns false.
\end{definition}

\begin{definition}[One evaluator step]
\lean{BoolCircuit.CircuitEval.absStep}
\label{BoolCircuit.CircuitEval.absStep}
\leanok
\uses{BoolCircuit.CircuitEval.AS, BoolCircuit.CircuitEval.absOf, BoolCircuit.CircuitEval.lact, BoolCircuit.CircuitEval.markCell}
For a configuration with phase $\phi$, computed Boolean values $v$, marked cells $M$,
and head position $j$, and an input symbol that may be absent, first apply the
evaluator's transition rule to $\phi$, the symbol, the value at position $j$ if present,
and the mark-tape lookup at $j$. Then apply the resulting action to the configuration
using the abstract update rule; the result is either a new configuration or an output
Boolean.
\end{definition}

\begin{definition}[Evaluator run]
\lean{BoolCircuit.CircuitEval.absRun}
\label{BoolCircuit.CircuitEval.absRun}
\leanok
\uses{BoolCircuit.CircuitEval.AS, BoolCircuit.CircuitEval.absStep}
A run on a finite Boolean input processes its symbols from left to right, starting from
a given evaluator configuration. At each symbol it takes one evaluator step; if that
step returns a new configuration, the run continues on the remaining input, and if it
returns an output Boolean, that Boolean is the result. Once the input is exhausted, the
evaluator takes a step with an absent symbol; if that step returns a configuration, the
result is false, and otherwise its output Boolean is the result.
\end{definition}

\begin{definition}[Leading run of ones]
\lean{BoolCircuit.CircuitEval.leadOnes}
\label{BoolCircuit.CircuitEval.leadOnes}
\leanok
For a finite Boolean string, the number of leading ones is the length of its longest
initial segment consisting entirely of true entries.
\end{definition}

\begin{definition}[Verifier verdict]
\lean{BoolCircuit.CircuitEval.verdict}
\label{BoolCircuit.CircuitEval.verdict}
\leanok
\uses{BoolCircuit.CircuitEval.absRun, BoolCircuit.CircuitEval.leadOnes, Turing.pairDecode}
Given a Boolean flag $e$ and a finite Boolean string $z$, attempt to decode $z$ as a
pair $(c,x)$ of Boolean strings; if decoding fails, the verdict is false. Otherwise let
$r$ be the number of leading ones in $c$. The verdict is false unless $c$ contains a
false entry, $r$ is at most the length of $x$, and either $e$ is false or $r$ equals the
length of $x$. When these conditions hold, the verdict is the result of running the
evaluator on $c$, starting in the phase that skips $r$ input positions, with the first
$r$ entries of $x$ as its computed values, no marked cells, and head position zero.
\end{definition}
